\addcontentsline{toc}{chapter}{Resumen}

\chapter*{Resumen}

\vspace*{-8mm}

Se presenta una implementación de Redes Neuronales Informadas por Física con
activación sinusoidal (PINN-SIREN) para la resolución de la ecuación de
Helmholtz en un paradigma libre de malla, validada contra soluciones
analíticas en 1D y en 2D con campo eléctrico complejo.
La arquitectura SIREN propuesta ($5 \times 128$ neuronas, $\omega_0 = 1.0$)
alcanza errores $L^2$ de $0.006\%$ en el caso 1D y $0.171\%$ en el caso 2D
respecto a soluciones analíticas, varios órdenes de magnitud por debajo del
umbral ingenieril de $5\%$ adoptado en este trabajo.
Una estrategia de optimización bifásica Adam~+~L-BFGS y el muestreo por
hipercubo latino (LHS) son clave para la convergencia y la precisión.
La reproducibilidad del método fue verificada con tres semillas independientes,
obteniendo un error $L^2$ promedio de $0.155 \pm 0.020\%$.
La generación de patrones de speckle óptico con condición de frontera de fase
aleatoria constituye el siguiente paso planificado de esta investigación.

\bigskip

\noindent\textbf{Palabras clave:} Redes Neuronales Informadas por Física,
SIREN, ecuación de Helmholtz, speckle óptico, activación sinusoidal,
muestreo por hipercubo latino.
